\documentclass[11pt]{article}
\usepackage[margin=1in]{geometry}
\usepackage[T1]{fontenc}
\usepackage[utf8]{inputenc}
\usepackage{lmodern}
\usepackage{microtype}
\usepackage{booktabs}
\usepackage{longtable}
\usepackage{array}
\usepackage{url}
\usepackage[hidelinks]{hyperref}
\usepackage{xcolor}
\usepackage{enumitem}
\setlist{nosep}
\newcommand{\sha}[1]{\texttt{\small #1}}
\newcommand{\file}[1]{\texttt{\small #1}}
\newcolumntype{P}[1]{>{\raggedright\arraybackslash}p{#1}}

\title{The Receipts Loop:\\Crash-Only, Fail-Closed Orchestration of LLM Coding Agents in Aesop,\\with a Dated Record of Its Practices}
\author{Matt Culliton\\\small Independent; \url{https://github.com/matt82198/aesop}}
\date{October 2026 \quad (manuscript draft; evidence snapshot at commit \sha{d8135aba}, 2026-10-01)}

\begin{document}
\maketitle

\begin{abstract}
Aesop is an open repository (first public commit 2026-07-11) that runs fleets of large-language-model coding agents against a real codebase under a set of operating rules enforced in code rather than in prompts. This paper is a record, not a proposal. The practices are these: crash-only orchestration with disk-resident state; an isolated orchestrator that dispatches small stateless workers; single-writer control files; fail-closed gates; regression traps generated from an incident log; detection of ``fake-green'' checks that report success without executing; a fail-closed gate against ``zombie'' work items, which are items still open after the work they describe has shipped; merge completion established by querying the hosting platform for the pull request's merged state, since a tool's exit code is not accepted as proof; and a rule against fitting assertions to observed output. For each practice we give the commit hash and ISO date at which it entered the repository, the incident that motivated it where one is logged, and the public essay or page that described it. We also report the project's measurements (a dispatch-topology A/B that found a hierarchical specialist tier cost 4--4.3$\times$ for identical output; a 39-task judgment benchmark on which the smallest model tier scored 39/39; a 130-task, five-tier, three-repeat benchmark) and the project's own documented failures, including seven pre-push checks found to fail open and a cost kill-switch that was wired but inert for a period. We make no claim that these practices are novel in software engineering; several are named in the project's own documents as descendants of crash-only software, Erlang supervision and write-ahead logging. The loop at the core is Huntley's Ralph Wiggum loop \cite{huntley2025}, a deliberately dumb \texttt{while} loop that re-feeds an agent its plan from disk; Aesop leaves the loop dumb and adds the rule that an iteration counts only when it leaves receipts on disk that a gate can check. We call the resulting pattern the receipts loop. We claim only that this dated record exists and is reproducible from the repository's history.
\end{abstract}

\section{Scope and evidentiary standard}

Every factual statement in this paper about Aesop is tied to one of three sources: a commit in the public repository \url{https://github.com/matt82198/aesop}, identified by hash and author date in ISO 8601 as reported by \texttt{git log --date=iso}; a committed document in that repository, identified by path and, where useful, line; or a dated public essay or page whose date is taken from the project's own curated index. The compiled evidence tables that back this paper are committed alongside it (\file{docs/paper/EVIDENCE.md} and \file{docs/paper/evidence.json} on branch \file{docs/arxiv-aesop}), and every row in them was produced by a mining pass over the repository on 2026-10-05 against the main-branch commit \sha{d8135aba} (2026-10-01).

Three kinds of statement are deliberately absent. The first is priority over other work: a single repository's history cannot say who did any of these things first, and the project's own essays trace most of the ideas to systems literature decades older than language models. The second is adoption by others. The third is any citation of tooling that lives on the operator's machine outside the repository (an idle-tick scheduler, a usage watchdog, a per-user dispatch hook), even where the repository's documentation refers to it, because its history cannot be dated from the public record.

Where the evidence contradicts a claim made inside the repository, we report the contradiction. Two such cases appear in Section~\ref{sec:limits}.

\section{The system in one paragraph}

Aesop is an orchestration harness: a single long-lived orchestrator process decomposes work into tasks and dispatches each to a short-lived worker agent in its own git worktree. What comes back is verified against gates that run as code, then opened as a pull request that a human or an automated merge tool lands. The orchestrator is forbidden from reading source code; it reads a small set of control files. Workers are forbidden from merging. Rules that must hold are implemented as pre-commit and pre-push hooks, continuous-integration jobs, and tests that fail when a previously logged failure pattern re-emerges. The design premise, stated in the project's first essay on 2026-07-12, is that the agent's own report of its work is not evidence of the work, so the harness checks. It does not trust \cite{hypothesis}.

\paragraph{Repository facts.} First public commit \sha{ca6437de} on 2026-07-11T15:50:02--05:00, ``initial open-source release of the fable-fleet orchestration harness''; a second commit the same day (\sha{5edf7313}) ported machinery from a private predecessor, including a document on surviving an endpoint-security agent that deleted the working tree (\file{docs/av-resilience.md}). As of \sha{d8135aba} (2026-10-01) the main branch carries 2,157 commits; the GitHub API reports 720 merged pull requests, against 865 merge commits in the history (the two are not interchangeable; the API count is the one we cite). Nineteen release tags run from \texttt{v0.1.0-beta.1} (2026-07-11) to \texttt{v0.8.0} (2026-09-11). The license changed three times. MIT became PolyForm Strict on 2026-07-17 (\sha{e1557a62}); it returned to MIT on 2026-07-29 (\sha{d8b4d829}); it became PolyForm Noncommercial on 2026-09-11 (\sha{07732210}).

\section{The loop at the core}
\label{sec:loop}

The outermost mechanism in Aesop is not clever, and that is inherited. On 2025-07-14 Geoffrey Huntley published the Ralph Wiggum loop \cite{huntley2025}: in its purest form, \texttt{while :; do cat PROMPT.md | claude-code ; done}. The agent is run again and again against a plan file on disk; whatever it forgets between runs it re-reads; whatever it breaks, the next iteration meets as a fact in the repository. Huntley's own description is that the technique is ``deterministically bad in an undeterministic world,'' and that is the point: a loop with no memory and no judgment cannot be talked into believing its own progress, because it has none to believe. It solves the long-running-agent problem by refusing to have a long-running agent.

Aesop's loop is that loop with a scheduler and a wall around it. The implementation is \file{driver/wave\_loop.py} (first added 2026-07-21, \sha{d6479458}); the operator-facing skill that drives a wave of it, \file{skills/buildsystem/SKILL.md}, dates from 2026-07-17 (\sha{9d792479}). The 2026-09-19 essay states the lineage directly: the outermost thing in Aesop ``primes from files, fills a backlog, empties it, and wakes itself up again. That is Ralph with a scheduler attached'' \cite{nothingnew}. The loop body stays dumb. What Aesop adds is on the way out of each iteration: nothing an iteration did is counted unless it left something on disk that a gate can check without asking the agent. A pull request is counted when the hosting platform reports it merged, not when a tool exits zero (Section~\ref{sec:practices}); a test is counted when it has been watched to fail; a fix is counted when the incident that motivated it has a trap; a wave is counted when its checkpoint is committed. The project's own word for these artifacts is \emph{receipts}, and it has published them as such since 2026-07-29 (\file{docs/RECEIPTS.md}, \sha{f6f2075f}). Hence the name we give the pattern: a Ralph loop whose iterations must produce receipts, and whose every failure mode is a restart from disk rather than a repair in place. Everything in Section~\ref{sec:practices} is a receipt type or a gate that demands one.

\section{Practices, dated}
\label{sec:practices}

Each subsection names the practice and its enforcing mechanism, then gives the earliest in-repository evidence. Table~\ref{tab:provenance} collects the dates.

\subsection{Crash-only orchestration and disk-resident state}

The orchestrator treats its own context as a cache. The filesystem is the memory. Intent and next steps live in \file{STATE.md} (seeded 2026-07-12, \sha{df73dc1b}); an append-only \file{BUILDLOG.md} records checkpoints (first added 2026-07-25, \sha{9584b11a}); \file{docs/CHECKPOINTING.md} states the rule as ``Never invent state. Verify everything from disk \ldots\ Disk is the source of truth.'' A session can end at any instant and the next one resumes from these files. The project's design rationale and measured recovery behaviour are in \file{docs/WHY-CRASH-ONLY.md} and \file{docs/crash-only-whitepaper.md}, both added 2026-07-29 (\sha{aa875898}). The same discipline was applied to the model seam: the orchestrator and worker seats became hot-swappable on 2026-07-25 (\sha{ebaf5dd1}), documented in \file{docs/MICROKERNEL.md} (\sha{11ea50f3}). The motivating incident, an endpoint-security agent that repeatedly deleted the working tree, is described in the 2026-07-12 essay \cite{hypothesis}; the design term is taken from Candea and Fox \cite{candea2003}.

\subsection{Orchestrator isolation and flat dispatch}

Rule 4 of \file{docs/CARDINAL-RULES.md}, present at the initial commit, reads ``Orchestrator reads ONLY: cardinal rules, STATE.md, BUILDLOG.md, MEMORY.md, and short git one-liners. Dispatch Haiku for research.'' \file{docs/DISPATCH-MODEL.md}, same commit, states ``subagents ALWAYS Haiku unless scoped work genuinely exceeds its capability (rare).'' The consequence is that the orchestrator's context does not grow with the fleet; workers return briefs, never transcripts. The project tested the alternative before rejecting it: a spike on 2026-07-13 (\sha{0faf81ab}) built a hierarchical topology with an intermediate specialist tier, and two independent A/B measurements on 2026-07-14, published with raw data on 2026-07-29 (\sha{43f99de5}, \file{docs/ab-cost-dataset.md}), found the tier ``increases cost 4--4.3$\times$ while producing identical quality (100\% test pass, zero repair rounds).'' The flat design was kept. On 2026-09-11 the worker-side rules were consolidated into \file{LANE-CONTRACT.md} (\sha{8837d4e9}): ``lanes are stateless processes. A lane that fails is KILLED, never corrected.''

\subsection{Single-writer control files}

Rule 7 at the initial commit: ``Single-writer control files: MEMORY.md (keeper writes only), STATE.md (orchestrator writes only), BUILDLOG.md (append-only for everyone, never overwrite).'' \file{docs/GOVERNANCE.md} extends this with an inbox pattern: processes that are not the designated writer append requests to an inbox file; they never edit control files directly. The practical effect is that concurrent agents cannot corrupt the orchestrator's view of the world; they can only queue for it.

\subsection{Fail-closed gates, and an audit of which ones were not}

The secret scanner that gates every push was ported on 2026-07-11 (\sha{a27c4f8a}). Its history is a record of closing holes: the push gate was found inert and restored on 2026-07-15 (\sha{c8b3c257}); worktree and blob bypasses were closed on 2026-07-16 (\sha{de1ddac7}); and on 2026-07-21 it was changed to fail closed on file or git read errors (\sha{dc765865}), a change the commit itself classifies as a P1 security fix. Further gates followed: a self-updating test-suite-count gate on 2026-07-26 (\sha{05d3e798}), a gate that blocks pushes which change behaviour without updating the governing \file{CLAUDE.md} on 2026-07-30 (\sha{92bad310}), and a generated tool index with a check mode on 2026-08-03 (\sha{c0d81133}). A kill switch and a cost ceiling were added on 2026-07-16 (\sha{c03e1419}).

The project then audited its own gates. \file{docs/GATES-FIRED.md} (2026-07-31, \sha{3f0aebab}) lists, with citations, which guardrails had actually fired, and reports as its critical finding that ``Seven pre-push checks in hooks/pre-push-policy.sh deliberately fail-open when their tool file is missing.'' The 0.7.1 changelog entry of the same day (\sha{d425709c}) states that ``The cost kill-switch is inert'': the ceiling was wired to four abort checkpoints but the wave loop never wrote the ledger it evaluated, so it could never fire. A later changelog entry claims the kill switch was wired end-to-end; the evidence pass did not isolate the commit that did so, so we report it as a claim the pass did not verify.

\subsection{Fake-green: checks that report success without running}
\label{sec:fakegreen}

The incident log defines a class \texttt{fake-green} as ``Tests reported green but never ran or skipped real validation'' and records two occurrences. The documented instance is a continuous-integration job of browser proofs that passed while executing no specifications; the fix, ``actually execute playwright specs + minimal dashboard smoke,'' landed on 2026-07-28 (\sha{8873971c}). The next day a regression trap was committed: \texttt{TestFakeGreenTrap} in \file{tests/test\_traps.py} (2026-07-29, \sha{b796876c}) fails if the pattern re-emerges. A static linter for workflow files, \file{tools/ci\_workflow\_lint.py}, is described in the changelog as catching ``the green-means-never-ran class.'' The rule was later generalised in \file{LANE-CONTRACT.md} (2026-09-11): ``Green can mean never ran,'' with a required discipline the document calls \emph{watched-to-fail}: break the behaviour, run the test, confirm it goes red, revert, confirm green. The phenomenon was described publicly in the 2026-07-22 field note \emph{Green Is Not Correct} \cite{green}, six days before the browser-proofs fix. One correction to the repository's own record belongs here: the trap's docstring dates the triggering commit to 2026-07-13; git dates it 2026-07-28, and we use the git date.

\subsection{Zombie work items and the resurrection gate}

A ``zombie'' in the project's vocabulary is a tracker item that remains open after the work it describes has shipped. On 2026-07-29 \file{docs/RECEIPTS.md} (\sha{f6f2075f}) reported that ``15/19 active tracker items were already shipped (79\% zombie rate)'' and named the structural fix. That fix arrived the same day as \file{tools/tracker\_guard.py} (\sha{7c69752a}), ``append-only lane journal + zombie-resurrection fail-closed gate,'' followed by automatic closure on pull-request merge (Guardrail G1, \sha{47d1141c}; \file{tools/tracker\_autoclose.py}, \sha{bae1f9aa}) and a reconciliation tool on 2026-07-30 (\file{tools/tracker\_reconcile.py}, \sha{9bce28d7}). The 79\% figure was re-verified by a metrics gate on 2026-07-30 (\sha{2b475ea3}). Six guardrail gates, including the tracker guard, were wired as enforcement in continuous integration on 2026-07-29 (\sha{1599f42c}).

\subsection{The incident-to-trap pipeline}

\file{docs/INCIDENTS.md} tracks operational failures by class with detection, resolution and source reference; it was present by 2026-07-12. On 2026-07-29 two things were added: a generator that derives the failure taxonomy from the committed record (\file{tools/incident\_report.py}, \sha{0902d05d}) and the adversarial trap suite (\file{tests/test\_traps.py}, \sha{b796876c}), whose header states ``Each trap encodes a pattern found in docs/INCIDENTS.md and fails if that pattern re-emerges.'' The suite covers the classes fake-green, test-pollution, gate-activation and doc-invented, and states explicitly which classes it cannot mechanise: conflict, flake, stall and ci-drift. At the evidence snapshot the log holds 47 incidents: stall 16, gate-activation 7, conflict 6, flake 6, test-pollution 6, ci-drift 3, fake-green 2, doc-invented 1. The policy that a red CI run goes to a resolver that must also build the missing gate is documented on 2026-07-30 (\sha{b1514f07}); the detector that implements it lives outside the repository and is not claimed here.

\subsection{Merge completion verified against platform state}

A deterministic serial merge train was added on 2026-07-29 (\file{tools/merge\_train.py}, \sha{3644fc99}), a batch auto-merge tool on 2026-07-30 (\file{tools/auto\_merge.py}, \sha{1de9088e}), and a shell-injection hole in the latter was closed the same day (\sha{46221764}). The rule these tools serve was written down in \file{LANE-CONTRACT.md} on 2026-09-11: ``MERGED-state proof via GitHub API, not by assumption \ldots\ An HTTP 200 response \ldots\ is NOT proof,'' and lanes ``NEVER merge.'' The document narrates the incident behind the rule: a merge train was handed a branch name copied from a worker's report, merged a stale branch, every gate passed, the bar was green, and the feature was never on main.

\subsection{Never fit green}

\file{LANE-CONTRACT.md} \S3 (2026-09-11): ``Never relax an assertion, lower a floor/ratchet, delete a suite, add a skip, or retune a constant to make a tree green. If an expectation is stale, the OWNING LANE derives the correct value \ldots\ A ratchet may be RAISED by hand with a measurement and a stated reason.'' This is the only practice in this paper whose in-repository statement postdates its use by the operator; earlier instances exist only in records outside the repository, so 2026-09-11 is the date we give.

\subsection{Watchdogs and stall detection}

A watchdog daemon script was present at the initial commit (\file{daemons/run-watchdog.sh}). Silent-hang detection for agents was added on 2026-07-13 (\file{tools/stall\_check.py}, \sha{17010688}) and hardened the same day against dangling-symlink and path-traversal inputs (\sha{8c3a7952}, \sha{504d864a}). The kill switch and cost ceiling of 2026-07-16 are the fleet-level brake; their documented inert period is covered above.

\subsection{Writing-side gates}

On 2026-10-01 three tools were added within one minute of each other: a detector of machine-writing patterns with a scored lint (\file{tools/humanize\_lint.py}, \sha{73643e36}), an append-only ledger of before/after edits with the rule each edit cleared (\file{tools/humanize\_ledger.py}, \sha{58c52fa6}), and a style-profile builder with a divergence report (\file{tools/humanize\_voice.py}, \sha{9200ba7f}). They apply the same incident-to-gate discipline to prose the agents produce. They postdate the \texttt{v0.8.0} tag and are unreleased at the snapshot.

\section{Measurements}
\label{sec:measure}

The following are the project's own measurements, with the committed source for each. We report them as the repository does and do not extend them.

\begin{itemize}
\item \textbf{Dispatch topology.} Two independent A/B measurements on 2026-07-14 (\file{docs/ab-cost-dataset.md}, published 2026-07-29, \sha{43f99de5}): a hierarchical topology with an intermediate specialist tier cost 4--4.3$\times$ the flat topology for identical quality on the same fixtures. The hierarchical design was cancelled.
\item \textbf{Judgment benchmark.} The v3 run of 2026-07-17 (\sha{21970ab1}; \file{bench/results/2026-07-17-judgment-v3-haiku-sonnet-opus.md}): on 39 held-out judgment tasks, the smallest tier scored 39/39, the middle tier 39/39, the largest 38/39. The project treats this as sufficiency for the work, not parity; a pre-registered ceiling rule applies.
\item \textbf{Frontier v4.} Published 2026-07-27 (\sha{10dbae11}; \file{bench/results/frontier-v4-2026-07-27.md}): 130 tasks, five model tiers, three repeats, 1,950 tuples over direct API transports, 1,742 scorable. Pairwise rows are reported with 90\% confidence intervals; the committed table lists, for example, haiku-4.5 versus sonnet-5 at $-7.69$ points with interval $[-11.70, -3.68]$ and sonnet-5 versus gpt-4o-mini at $+13.85$ with interval $[9.53, 18.16]$, both marked indeterminate under the pre-registered equivalence margin.
\item \textbf{Refusal ladder and seated A/B.} A ten-model refusal ladder on a blind 16-item corpus was published on 2026-07-23 (\sha{1ba5a2a6}); a seated A/B over the completed model seam on 2026-07-24 (\sha{fd8da971}) reports that supplying real repository context changed the verdict on one hard item for two of three tested models. Both are small-N and are labelled as such in their source files.
\end{itemize}

\section{Limitations and counter-evidence}
\label{sec:limits}

\paragraph{Single operator, single repository.} Everything above was produced by one person and one fleet on one codebase, the harness's own. Nothing here has been replicated by an independent party, and the paper makes no claim that it would be.

\paragraph{The gates were not all closed.} Seven pre-push checks were found to fail open on 2026-07-31; the cost kill-switch was inert for a period the project itself documented. We include these because a paper about gates that did not report the gates that did not fire would be an instance of the fake-green class it describes.

\paragraph{Internal record errors.} A test docstring misdates its own triggering commit by fifteen days (Section~\ref{sec:fakegreen}). The merge-commit count (865) and the merged-pull-request count (720) differ and are not interchangeable. The exact commit that resolved the inert kill-switch was not isolated.

\paragraph{Machine-local tooling.} Several mechanisms the repository's documentation depends on (an idle-tick scheduler, a usage watchdog, a red-CI detector, a dispatch hook that pins worker models) live outside the repository and are excluded from every claim here.

\paragraph{No priority, no adoption.} The loop is Huntley's \cite{huntley2025}. The practices around it are, by the project's own account \cite{nothingnew}, applications of crash-only software \cite{candea2003}, Erlang-style supervision \cite{armstrong2003}, write-ahead logging \cite{mohan1992}, leases \cite{gray1989} and fail-safe defaults \cite{saltzer1975} to a component that can fail and can also report success falsely. What this paper contributes is the dated record of those applications, in the form the repository can support.

\section{Reproducing this record}

Clone the repository, check out \sha{d8135aba}, and run \texttt{git log --date=iso --reverse -S<term>} for any term in Table~\ref{tab:provenance}; the dates in this paper are the author dates of the first matching commits. The evidence tables (\file{docs/paper/EVIDENCE.md}, \file{docs/paper/evidence.json}) record every query that produced a row, and the ten ambiguities the mining pass flagged are listed at the end of that file.

\appendix
\section{Provenance table}

\begin{longtable}{P{3.6cm}P{2.3cm}P{2.3cm}P{6.2cm}}
\caption{First in-repository evidence for each practice. Dates are commit author dates (ISO, local offset $-$05:00). Public descriptions are dated from the project's curated index.}\label{tab:provenance}\\
\toprule
Practice & Date & Commit & Evidence and public description \\
\midrule
\endfirsthead
\toprule
Practice & Date & Commit & Evidence and public description \\
\midrule
\endhead
Orchestrator isolation (reads only control files) & 2026-07-11 & \sha{ca6437de} & \file{docs/CARDINAL-RULES.md} rule 4; essay 2026-07-24 \cite{unix} \\
Wave skill (the operator's loop entry point) & 2026-07-17 & \sha{9d792479} & \file{skills/buildsystem/SKILL.md} \\
Wave loop implementation (Ralph loop with gates) & 2026-07-21 & \sha{d6479458} & \file{driver/wave\_loop.py}; Ralph loop published 2025-07-14 \cite{huntley2025} \\
Flat dispatch, smallest-model workers & 2026-07-11 & \sha{ca6437de} & \file{docs/DISPATCH-MODEL.md}; A/B 2026-07-14; essay 2026-07-18 \cite{haiku} \\
Single-writer control files & 2026-07-11 & \sha{ca6437de} & \file{docs/CARDINAL-RULES.md} rule 7; \file{docs/GOVERNANCE.md} \\
Secret-scan push gate & 2026-07-11 & \sha{a27c4f8a} & fail-closed on read errors 2026-07-21 (\sha{dc765865}) \\
Crash-only control files & 2026-07-12 & \sha{df73dc1b} & \file{STATE.md}; \file{BUILDLOG.md} 2026-07-25; whitepaper 2026-07-29 (\sha{aa875898}); essay 2026-07-12 \cite{hypothesis} \\
Incident log & by 2026-07-12 & --- & \file{docs/INCIDENTS.md}; 47 incidents at snapshot \\
Stall detection & 2026-07-13 & \sha{17010688} & \file{tools/stall\_check.py} \\
Hierarchical tier tested and rejected & 2026-07-13/14 & \sha{0faf81ab} & dataset published 2026-07-29 (\sha{43f99de5}) \\
Kill switch and cost ceiling & 2026-07-16 & \sha{c03e1419} & inert period documented 2026-07-31 (\sha{d425709c}) \\
Judgment benchmark 39/39 & 2026-07-17 & \sha{21970ab1} & \file{bench/results/2026-07-17-judgment-v3-haiku-sonnet-opus.md} \\
Fake-green: public description & 2026-07-22 & --- & field note \cite{green} \\
Hot-swappable seats & 2026-07-25 & \sha{ebaf5dd1} & \file{docs/MICROKERNEL.md}; artifact 2026-07-25 \cite{seats} \\
Test-suite-count gate & 2026-07-26 & \sha{05d3e798} & \file{tools/verify\_test\_suite\_count.py} \\
Frontier v4 benchmark & 2026-07-27 & \sha{10dbae11} & \file{bench/results/frontier-v4-2026-07-27.md}; artifact \cite{frontier} \\
Fake-green: incident fixed & 2026-07-28 & \sha{8873971c} & browser proofs executed; class \texttt{fake-green} in \file{docs/INCIDENTS.md} \\
Zombie-resurrection gate & 2026-07-29 & \sha{7c69752a} & \file{tools/tracker\_guard.py}; 79\% finding \file{docs/RECEIPTS.md} (\sha{f6f2075f}) \\
Incident-to-trap suite & 2026-07-29 & \sha{b796876c} & \file{tests/test\_traps.py}; generator \sha{0902d05d} \\
Serial merge train & 2026-07-29 & \sha{3644fc99} & \file{tools/merge\_train.py}; auto-merge 2026-07-30 (\sha{1de9088e}) \\
Gates-fired audit (seven fail-open) & 2026-07-31 & \sha{3f0aebab} & \file{docs/GATES-FIRED.md} \\
Generated tool index with check & 2026-08-03 & \sha{c0d81133} & \file{tools/gen\_tool\_index.py} \\
Lane contract: MERGED-state proof, watched-to-fail, never fit green & 2026-09-11 & \sha{8837d4e9} & \file{LANE-CONTRACT.md} \S3, \S5, \S6, \S7 \\
Writing-side gates (humanize) & 2026-10-01 & \sha{73643e36} & \file{tools/humanize\_lint.py}, \file{humanize\_ledger.py}, \file{humanize\_voice.py} \\
\bottomrule
\end{longtable}

\begin{thebibliography}{9}
\bibitem{huntley2025} G.~Huntley. Ralph Wiggum as a ``software engineer''. ghuntley.com, 2025-07-14. \url{https://ghuntley.com/ralph/}
\bibitem{hypothesis} M.~Culliton. The Aesop Hypothesis: AI agents that survive because they're designed to fail. Medium, 2026-07-12. \url{https://medium.com/@matt82198/the-aesop-hypothesis-ai-agents-that-survive-because-theyre-designed-to-fail-de5f033369d4}
\bibitem{haiku} M.~Culliton. The Haiku Wager: autonomous dev at 1/3 the cost. Medium, 2026-07-18. \url{https://medium.com/@matt82198/the-haiku-wager-autonomous-dev-at-1-3-the-cost-93c0ecf4c0ec}
\bibitem{green} Aesop orchestrator (published by M.~Culliton). Green Is Not Correct: field notes from an AI orchestrator. Medium, 2026-07-22. \url{https://medium.com/@matt82198/green-is-not-correct-field-notes-from-an-ai-orchestrator-end-of-a-long-shift-099ae33a724a}
\bibitem{unix} M.~Culliton. We don't need smarter LLMs. We need Unix for agents. Medium, 2026-07-24. \url{https://medium.com/@matt82198/we-dont-need-smarter-llms-we-need-unix-for-agents-d3c10d9f6f8f}
\bibitem{nothingnew} M.~Culliton. Nothing Here Is New, Except Chatter. Medium, 2026-09-19. \url{https://medium.com/@matt82198/nothing-here-is-new-except-chatter-d4404b0e106f}
\bibitem{seats} aesop 0.4.0: Two Swappable Seats. Claude artifact, 2026-07-25. \url{https://claude.ai/artifact/BTwxP8mxbHtJFsT7Hm8bPU}
\bibitem{frontier} Frontier v4: Final Results. Claude artifact, 2026-07-27. \url{https://claude.ai/artifact/24pkx9kHFXFnGE5EH9LYEC}
\bibitem{candea2003} G.~Candea and A.~Fox. Crash-only software. In \emph{HotOS IX}, 2003.
\bibitem{armstrong2003} J.~Armstrong. \emph{Making reliable distributed systems in the presence of software errors}. PhD thesis, KTH, 2003.
\bibitem{mohan1992} C.~Mohan et al. ARIES: a transaction recovery method supporting fine-granularity locking and partial rollbacks using write-ahead logging. \emph{ACM TODS}, 17(1), 1992.
\bibitem{gray1989} C.~Gray and D.~Cheriton. Leases: an efficient fault-tolerant mechanism for distributed file cache consistency. In \emph{SOSP}, 1989.
\bibitem{saltzer1975} J.~H.~Saltzer and M.~D.~Schroeder. The protection of information in computer systems. \emph{Proceedings of the IEEE}, 63(9), 1975.
\end{thebibliography}

\end{document}
